\section{Introduction}
\label{sec:intro}

Evidence grounding is a standard recipe for making language-model fact verification more trustworthy.
If a model is shown the same designated evidence that a benchmark treats as sufficient, aggregate accuracy should rise, and remaining errors should look like residual retrieval or coverage problems rather than failures to use gold evidence.
That picture is incomplete.
Aggregate improvements can hide a second phenomenon: claims the model already answered correctly without evidence, then answered incorrectly after the designated evidence was added.

This paper studies that phenomenon on FEVER \citep{thorne-etal-2018-fever}.
We evaluate GPT-5.4 and GPT-5.4-mini on a balanced 10{,}000-claim Supported/Refuted sample under two paired conditions: claim only, and claim plus FEVER-designated gold evidence.
Evidence helps in aggregate (Table~\ref{tab:accuracy}).
It also induces 226 unique correct-to-wrong regressions.
Those regressions are rare relative to the 10{,}000-claim sample, but they are not unstructured noise.
They concentrate on claim--evidence pairs that local NLI models treat as weakly warranted, and they are only partly shared across the two GPT models.

Counting regressions is not enough.
The same behavioral pattern can arise from at least three diagnostic situations: the model fails to use evidence that a later adjudicator finds decisive; the evidence representation (sentence text versus titles versus structured provenance) changes the warrant; or the claim--evidence pair remains underdetermined after full designated evidence is shown.
We therefore run blinded progressive-disclosure silver adjudication.
Five isolated judges see sentence evidence (Stage~A), then page titles (Stage~B), then reconstructed structured FEVER evidence (Stage~C).
A rule-based taxonomy maps consensus transitions onto three broad diagnostic categories: evidence-utilization failure, representation-sensitive behavior, and residual ambiguity.
The primary panel uses one locked Cursor/Grok model.
Supporting checks include a GLM-5.3 Stage~A replication and a Claude Stage~C test on 231 Grok leftovers.
The decisive robustness experiment is a complete Claude A$\rightarrow$B$\rightarrow$C panel on all 226 regressions.

The resulting claim is two-part.
First, designated gold evidence improves overall accuracy while creating a measurable set of localized regressions whose diagnostic picture is mixed rather than single-cause.
Second, that mixed picture is qualitatively stable across judge families, but exact category prevalences are not.
Automated diagnostic decompositions therefore carry measurement uncertainty of their own.

\paragraph{Contributions.}
We (i)~quantify paired evidence gains and evidence-induced regressions on a verified 10{,}000-claim FEVER experiment;
(ii)~characterize regressions with confirmatory statistics, two local NLI diagnostics, and a 1{,}061-claim diagnostic cohort;
(iii)~decompose the 226 unique regressions with blinded progressive-disclosure silver adjudication; and
(iv)~show, with a 3{,}390-judgment Claude replication, which parts of that decomposition survive a change of adjudicating family.

We ask four research questions, restated in Section~\ref{sec:rqs}.
RQ1 asks whether designated gold evidence improves aggregate accuracy.
RQ2 asks how frequent and cross-model-consistent the regressions are.
RQ3 asks which diagnostic failure modes appear under progressive disclosure.
RQ4 asks which of those conclusions survive a change in judge family.
